\section{Introduction}
\label{sec:intro}

Every system that stores, transmits, or processes information eventually
has to answer a question it almost never states out loud: what, exactly,
is the ``data'' it claims to be operating on? Database textbooks call data
``raw facts,'' not yet processed into meaning~\autocite{rob2013};
information-science texts describe it as a coherent organization of such
facts~\autocite{ratzan2004}; UNESCO adds that data must be ``expressed in
a standardized way, suitable for being communicated, interpreted, or
processed''~\autocite{unesco2008}; general dictionaries fall back on data
as ``facts or things known, used as the basis for inference''~\autocite{websters1934}.
None of these is wrong, exactly. None of them is formal enough to answer
the question a type system or a wire protocol actually needs answered:
what structure does a collection of facts have to have before it can be
indexed, copied, compared, or serialized at all?

That gap does not stay theoretical for long. It shows up in production
systems as a small, recognizable family of defects:

\begin{description}
  \item[Schema drift.]
    Two components silently disagree about what an index or field name
    denotes, even while operating on byte-identical data.
  \item[Type confusion.]
    A fixed byte sequence gets reinterpreted under a different type
    assumption than the one it was written under---the structural pattern
    behind an entire class of memory-safety and deserialization
    vulnerabilities~\autocite{cwe843}.
  \item[Unsound equality.]
    Two values are treated as equal because they share a representation,
    while quietly differing in the value set the comparison assumed.
  \item[Serialization ambiguity.]
    A wire format under-specifies its index set or value set just enough
    that two conformant implementations disagree about what the same bytes
    mean.
\end{description}

These four look unrelated on the surface---a caching bug has nothing to do
with a deserialization exploit, on the face of it---but they share a
single root. In each case, some assumption about the \emph{domain}, the
\emph{codomain}, or the \emph{mapping} underlying a piece of data was
never written down, and so it could never be checked. The rest of this
paper is an attempt to write that assumption down once, precisely enough
that it stops being implicit.

\subsection{Contribution}

Concretely, this paper:
\begin{enumerate}
  \item States the function-theoretic definition of data formally, with
        explicit well-definedness conditions
        (Section~\ref{sec:definition}), and works out its relationship to
        sequences, tuples, and multisets along the way.
  \item Places that definition against three existing frameworks
        (Section~\ref{sec:related}): Floridi's General Definition of
        Information, Zins' faceted classification, and the DIKW hierarchy.
  \item Derives four corollaries (Section~\ref{sec:implications}), each
        tying a specific structural violation to one of the four defect
        classes above, with a worked example for each.
  \item Discusses where the definition strains---streaming data,
        genuinely unstructured text---and what it does not yet claim to
        do (Section~\ref{sec:discussion}).
\end{enumerate}

\subsection{Scope}

One thing this paper deliberately does not try to do: it is not a
philosophical account of the full data--information--knowledge
progression, and it does not take a side in the older debate over whether
data are theory-laden or simply ``given'' in some stronger epistemic
sense~\autocite{floridi2005}. The question here is narrower and more
mechanical---what structure does data need before a computer can index,
copy, or compare it---and the paper stays inside that boundary on purpose.